% Cote 119: offprints, reports, courses, CV and letters (typed and printed part): archive cote 119, site page 216 (PDF page 217).
% typescript. Draft compared: none.
% AI-transcribed by Claude Opus 5.5 (claude-opus-5-5, adaptive thinking, effort high), Message Batches API, batch r119_full (msgbatch_01DwZNm5Y7YsKBWyAN4sFs8w), 2026-11-13; not yet checked by eye.
\scanpage{119/216}
\mspage{. 3 .}

l'esprit de la théorie des n-catégories, particulièrement
bien adaptée à l'interprétation géométrique des invariants
cohomologiques (cf. le livre cité de J. Giraud et
le travail de Mme M. Raynaud [$*$]).

\uline{Références} : J.L. Verdier, P. Deligne, D. Quillen,
P. Gabriel.

3. \uline{Topologie} ([16, SGA 4],[9])

Jusqu'à présent, c'est surtout le K-invariant
des espaces topologiques que j'avais introduit à l'occasion
de mes recherches sur le théorème de Riemann-Roch en
géométrie algébrique, qui a connu la fortune la plus
brillante, étant le point de départ de très nombreuses
recherches en topologie homotopique et topologie
différentielle. De nombreuses constructions que j'avais
introduites pour les besoins de la démonstration algébrique
du théorème de Riemann-Roch (telles les opérations $\lambda_{i}$
et leurs liens avec les opérations du groupe symétrique)
sont devenues pratique courante non seulement en géométrie
algébrique et en algèbre, mais également en topologie
et en théorie des nombres, notamment dans les travaux
de mathématiciens comme Atiyah, Hirzebruch, Adams, Quillen,
Bass, Tate, Milnor, Karoubi, Shih, etc...

Plus fondamental me semble néanmoins l'élar-%
gissement de la topologie générale, dans l'esprit de
la théorie des faisceaux (développée initialement par
J. Leray), contenu dans le point de vue des topos
[16, SGA 4]. J'ai introduit ces topos à partir de 1958
en partant du besoin de définir une cohomologie $\ell$-adique
des variétés algébriques (plus généralement, des schémas),
qui convienne à l'interprétation cohomologique des

\margin{216}
